\documentclass[fleqn]{llncs}

%include agda.fmt

%subst keyword a = "\keyword{" a "}"
%subst conid a = "\identifier{" a "}"
%subst varid a = "\identifier{" a "}"
%format Set = "\constructor{Set}"
%format Set1 = "\constructor{Set_1}"
%format MAPSTO = "\mathrel{\mapsto}"
%format Maybe = "\constructor{Maybe}"
%format nothing = "\constructor{nothing}"
%format just = "\constructor{just}"
%format Bool = "\constructor{Bool}"
%format true = "\constructor{true}"
%format false = "\constructor{false}"
%format Par = "\constructor{Par}"
%format return = "\constructor{return}"
%format _>>=_ = _ "{" >>= "}\kern-\underscorespace" _
%format fail = "\constructor{fail}"
%format catch = "\constructor{catch}"
%format assert = "\constructor{assert}"
%link assert assert_then_
%format then = "\constructor{then}"
%link then assert_then_
%format assert_then_ = assert _ then "\kern-\underscorespace" _
%link >>= _>>=_
%format (HIDE(x)) = "{}^{\color{gray}" x "}"

\newlength{\underscorespace}
\setlength{\underscorespace}{.8pt}

\newcommand{\keyword}[1]{\mathbf{#1}}
\newcommand{\identifier}[1]{\mathit{#1}}
\newcommand{\constructor}[1]{\mathsf{#1}}

\usepackage[hidelinks]{hyperref}
\renewcommand{\chapterautorefname}{Chapter}
\renewcommand{\sectionautorefname}{Section}
\renewcommand{\subsectionautorefname}{Section}
\renewcommand{\figureautorefname}{Figure}

\usepackage[color=yellow,textsize=footnotesize]{todonotes}

\usepackage{amsmath}

\newcommand{\reason}[1]{\quad\{\,\text{#1}\,\}}

\newcommand{\BiGUL}{\textsc{BiGUL}}
\newcommand{\PutGet}{\textsc{PutGet}}
\newcommand{\GetPut}{\textsc{GetPut}}

\newtheorem{Definition}{definition}

\begin{document}

\title{A Minimalist Datatype-Generic Putback-Based Bidirectional Programming Language}
\author{Hsiang-Shang Ko\kern.05em\inst{2} \and Tao Zan\kern.05em\inst{1} \and Zhenjiang Hu\kern.05em\inst{1,\,2}}
\institute{%
SOKENDAI (The Graduate University for Advanced Studies), Japan
\and
National Institute of Informatics, Japan}
\maketitle

\section{Introduction}
\label{sec:introduction}

\cite{Foster-lenses}

\section{Decidable partiality}
\label{sec:partiality}

By ``decidable partiality'' we mean that the transformations are actually total functions which wrap their results in the standard |Maybe| datatype:
\begin{code}
data ^^^Maybe (A : Set) : Set where
  ^^^just     : A -> Maybe A
  ^^^nothing  : Maybe A
\end{code}
This is different from the usual partial functions as formalised in terms of complete partial orders:
In \autoref{sec:BiGUL} we will need to define programs that produce some result or fail respectively when their input fails or produces some result, but this is not possible in the setting of complete partial orders, since such functions violate monotonicity, a necessary condition for computability.

\begin{code}
data ^^^Par : Set -> Set1 where
  ^^^return        : (HIDE({A : Set} ->)) A -> Par A
  ^^^_>>=_         : (HIDE({A B : Set} ->)) Par A -> (A -> Par B) -> Par B
  ^^^fail          : (HIDE({A : Set} ->)) Par A
  ^^^catch         : (HIDE({A B : Set} →)) Par A → (A → Par B) → Par B → Par B
  ^^^assert_then_  : (HIDE({A : Set} →)) Bool → Par A → Par A
\end{code}

\begin{code}
^^^runPar : {A : Set} -> Par A -> Maybe A
runPar (return x               )  = just x
runPar (mx >>= f               )  with runPar mx
runPar (mx >>= f               )  | just x   = runPar (f x)
runPar (mx >>= f               )  | nothing  = nothing
runPar fail                       = nothing
runPar (catch mx f my          )  with runPar mx
runPar (catch mx f my          )  | just x   = runPar (f x)
runPar (catch mx f my          )  | nothing  = runPar my
runPar (assert true   then mx  )  = runPar mx
runPar (assert false  then mx  )  = nothing
\end{code}

\begin{figure}
\begin{code}
mutual

  data CompSeq : {A : Set} → Par A → A → Set₁ where
    ret          :  {A : Set} {x x' : A} → x ≡ x' → CompSeq (return x) x'
    bind         :  {A B : Set} {x : A} {mx : Par A} {f : A → Par B} {y : B} →
                    CompSeq mx x → CompSeq (f x) y → CompSeq (mx >>= f) y
    catch-fst    :  {A B : Set} {mx : Par A} {f : A → Par B} {my : Par B} {x : A} {z : B} →
                    CompSeq mx x → CompSeq (f x) z → CompSeq (catch mx f my) z
    catch-snd    :  {A B : Set} {mx : Par A} {f : A → Par B} {my : Par B} {z : B} →
                    FailedCompSeq mx → CompSeq my z → CompSeq (catch mx f my) z
    assert-then  :  {A : Set} {b : Bool} {mx : Par A} {x : A} → b ≡ true → CompSeq mx x → CompSeq (assert b then mx) x

  data FailedCompSeq : {A : Set} → Par A → Set₁ where
    bind-fst   : {A B : Set} {mx : Par A} {f : A → Par B} → FailedCompSeq mx → FailedCompSeq (mx >>= f)
    bind-snd   : {A B : Set} {mx : Par A} {x : A} {f : A → Par B} →
                 CompSeq mx x → FailedCompSeq (f x) → FailedCompSeq (mx >>= f)
    fail       : {A : Set} → FailedCompSeq (fail {A})
    catch-fst  : {A B : Set} {mx : Par A} {f : A → Par B} {my : Par B} →
                 FailedCompSeq mx → FailedCompSeq my → FailedCompSeq (catch mx f my)
    catch-snd  : {A B : Set} {mx : Par A} {f : A → Par B} {my : Par B} {x : A} →
                 CompSeq mx x → FailedCompSeq (f x) → FailedCompSeq (catch mx f my)
    assert-fst : {A : Set} {mx : Par A} {b : Bool} → b ≡ false → FailedCompSeq (assert b then mx)
    assert-snd : {A : Set} {mx : Par A} {b : Bool} → b ≡ true → FailedCompSeq mx → FailedCompSeq (assert b then mx)
\end{code}
\caption{Definitions of |CompSeq| and |FailedCompSeq|.}
\label{fig:CompSeq}
\end{figure}

\section{Putback-based bidirectional transformations}
\label{sec:BX}

\begin{align*}
           & |get s MAPSTO v| \\
\Rightarrow& \reason{\GetPut\ for |get|} \\
           & |put s v MAPSTO s| \\
\Rightarrow& \reason{\PutGet\ for |get'|} \\
           & |get' s MAPSTO v|
\end{align*}

\section{\BiGUL}
\label{sec:BiGUL}

\section{Reduction of common features}
\label{sec:reduction}

\section{Discussion}
\label{sec:discussion}

\todo[inline]{Haskell implementation}

\bibliographystyle{plain}
\bibliography{/Users/joshko/Documents/bib}

\end{document}